\chapter{Conclusions}

This project built and evaluated a complete data-management pipeline for comparing
restaurant ratings across Google Maps, Tripadvisor, and TheFork in the Milan area.
Starting from three heterogeneous sources with independent identifiers, formats, and
levels of quality, the work produced a single integrated database suitable for
consistency analysis and analytical querying.

\section{Summary of the Work}

The pipeline follows an ELT design. Raw data are acquired through the Google Places
API and two dedicated scrapers, loaded into MongoDB as an immutable system of record,
and then cleaned into per-source collections. Schema matching aligns the three
vocabularies and classifies their representation conflicts, after which a
Google-anchored entity-resolution stage links records through blocking, similarity
scoring, and calibrated thresholds, with dedicated hardening for chain and franchise
venues and an optional large-language-model adjudication of uncertain pairs. The
selected links are materialized into a Google-seeded integrated collection that keeps
each platform's ratings side by side, reconciles cuisine into a canonical taxonomy,
and resolves price by majority-vote consensus. Finally, the cleaned and integrated
collections are projected into a flat ClickHouse layer that supports the
research-question queries.

Quality was measured at two distinct points. Before integration, a data-quality
assessment profiled each raw source along standard data-quality dimensions. After
integration, the assessment targeted the integration process itself, evaluating the
entity-resolution classifier, one-to-one link survival, and Tripadvisor geocoding
against a manually labelled gold standard.

\section{Main Findings}

The integrated dataset contains 10,054 restaurants. Ratings are generally consistent
across platforms, but the agreement is not perfect and its structure is informative.
Tripadvisor is systematically more severe than Google and TheFork, which are closely
aligned with each other (RQ1). Rating level and cross-platform agreement vary by
cuisine, with mainstream categories such as pizza and fast food being lower-rated and
more contested than specialist cuisines (RQ2). Data completeness is spatially
structured, with the richest metadata concentrated in the historic centre and the
Navigli and Isola belt and decreasing towards the periphery (RQ3).

\section{Limitations}

Several limitations follow from the design. The dataset is Google-seeded, so a
restaurant absent from Google cannot enter the integrated collection, and TheFork's
small footprint limits three-platform coverage. Tripadvisor coordinates are obtained
through geocoding rather than measured, which introduces a small spatial error.
Entity resolution is deliberately conservative: many true matches are left as
uncertain rather than forced, and one-to-one link selection reduces the survival of
some human-confirmed matches. Finally, the data represent a single point in time and
do not capture how ratings evolve.

\section{Future Work}

The integrated dataset opens several extensions. A first group of additional research
questions --- the effect of review volume on rating stability, the role of price level,
the impact of multi-platform presence on visibility, and the relationship between photo
activity and engagement --- is being carried forward as part of the data-visualization
module of the course. Beyond these, the most natural extension is to exploit the
review text that is already collected but not yet analysed, applying language models
to extract sentiment and topics and to explain rating differences qualitatively.
Periodic re-acquisition would add a temporal dimension and allow the study of rating
dynamics, while broadening the geographic scope beyond Milan would test whether the
observed platform biases generalize to other markets.
